%% kapa: exam=1 round=28 year=2017 period=1 form=A subject=real_estate subject_ko=부동산학원론 seq=12 number=92
%% source: Q-Net 감정평가사 기출문제 (cst00309) artlSeq=5199603 file=extracted/1008/5199603/2017년 제28회 감정평가사 1차시험 문제지/1교시 A형.pdf
%% conversion: kordoc 4.24.51 -> tools/gen1.py (auto); review=proofread-sonnet; flags=none
\begin{question}{92}
부동산투자시 위험과 수익과의 관계에 관한 설명으로 옳은 것을 모두 고른 것은?
\begin{framedtext}
ㄱ. 위험회피형 투자자는 위험 증가에 따른 보상으로 높은 기대수익률을 요구한다. \par
ㄴ. 위험과 수익과의 상쇄관계는 위험이 크면 클수록 요구하는 수익률이 작아지는 것을 의미한다. \par
ㄷ. 위험의 크기에 관계없이 기대수익률에만 의존해서 행동하는 투자유형을 위험선호형이라 한다. \par
ㄹ. 요구수익률은 무위험률과 위험할증률을 합산하여 계산해야 한다. \par
ㅁ. 평균-분산모형에서, 기대수익률이 같다면 위험이 작은 투자안을 선택하고, 위험이 같다면 기대수익률이 높은 투자안을 선택하는 투자안의 선택기준을 지배원리(dominance principle)라고 한다.
\end{framedtext}
\begin{choices}
  \item ㄱ, ㄴ
  \item ㄴ, ㄷ
  \item ㄱ, ㄹ, ㅁ
  \item ㄴ, ㄷ, ㅁ
  \item ㄷ, ㄹ, ㅁ
\end{choices}
\end{question}
\answer{3}
